\section{Cómo una carencia llega a propuesta}
\label{sec:esquema:metodo}\label{sec:propuesta:conformidad}\label{sec:evolucion:metodo}\label{sec:evolucion:clases}

Las ocho Datasheets del catálogo se escriben contra la guía de ACM2YAML v0.3, el documento que describe el esquema con ejemplos, sin inventar campos \citep{apisconsumptionmodel}.
Las dos formas que quedan fuera de la guía, la unidad \texttt{MBs} de los topes de tamaño y el registro \texttt{not\_found} de lo que el proveedor no publica, ya se describieron junto al consumo del catálogo por PRIME4API (Sección~\ref{sec:propuesta:prime4api}).
Escribir sin campos propios tiene una ventaja: el catálogo sirve aunque no se acepte ninguna extensión, porque PRIME4API lo lee tal cual, con los fallos que recoge la Tabla~\ref{tab:verificacion}.

Cuando un proveedor publica algo que el esquema no puede representar, no se añade a la Datasheet ningún campo propio.
El hecho se registra como «elemento no soportado»: un hecho del dominio que ACM2YAML v0.3 no puede expresar, como los seis volúmenes contratados de Postmark, o que expresa de una forma que PRIME4API, el motor de referencia, no consume.
Cada elemento lleva un número F-nn, el caso real que lo motiva, la consecuencia para el análisis y, cuando lo hay, un esbozo de resolución.
En el repositorio de datos estos registros se llaman «fricciones» (fichero \texttt{fricciones.md}), y de ahí viene la F del número.

Un elemento no soportado se promueve, es decir, pasa a candidato de extensión del esquema, por una de dos vías.
La primera es la recurrencia: el elemento reaparece en al menos dos APIs del catálogo.
Que Postmark venda seis volúmenes contratados para su cuota mensual podría ser una rareza suya, y verlo en un segundo proveedor indica que pertenece al dominio.

La segunda vía es la traza, que admite una sola API si se cumplen a la vez tres condiciones.
La condición L1 pide que el hecho esté citado con su literal o medido en un sondeo, con la evidencia archivada en el repositorio.
La condición L2 pide la consecuencia ejecutada sobre el motor: con la Datasheet sin manipular, PRIME4API devuelve un valor erróneo o deja sin responder lo que respondería si el esquema pudiera expresar el hecho.
Pide además un caso de control: la misma operación del motor responde bien cuando el dato cabe en el esquema.
La condición L3 pide que ningún mecanismo existente de la guía cubra la carencia y que no la explique un defecto ya conocido del motor, lo que descarta una mala lectura de la guía.

Quedar promovido no equivale a ser propuesta de extensión.
Solo lleva número de propuesta, P-nn, el elemento promovido cuya forma de extensión se ha comprobado contra los campos que define la guía.
La propuesta P-05 es la excepción: no sale de ningún elemento no soportado y trata el plan cuyo precio existe pero no se publica.

Un elemento que aparece en una sola API y no deja traza sobre el motor no se añade al esquema, aunque esté bien documentado.
Queda en el inventario sin promover, y su fila anota la condición que no cumple.
El criterio mide la generalidad dentro del catálogo, y no se ha comprobado si los once elementos no promovidos (Sección~\ref{sec:esquema:inventario}) aparecen en APIs de fuera del catálogo.

La recurrencia muestra que una carencia pertenece al dominio, pero no que la forma elegida para sortearla funcione.
Esa segunda comprobación exige una herramienta que lea la Datasheet y diga si la forma elegida se procesa.
Ninguno de los repositorios del ecosistema, ni el de la guía ni los de PRIME4API y SPHERE, incluye un esquema JSON, un analizador estático o una integración continua que compruebe un fichero \texttt{datasheet.yaml}.
PRIME4API es un motor de cálculo y no un validador: lee los campos del YAML sin comprobarlos antes, y sus esquemas describen las peticiones de su propia API REST \citep{prime4apirepo}.
El repositorio de la guía aporta la norma y los ejemplos, pero no un comprobador \citep{apisconsumptionmodel}.

Gavira Sánchez describe así el editor de Datasheets de SPHERE: «Cada cambio en el documento se valida contra la sintaxis del modelo» \citep{gavira2026tfm}.
En el repositorio de SPHERE, en su versión del 10 de junio de 2026, la página de subida exige la extensión \texttt{.yaml} o \texttt{.yml} y lee el YAML solo para sacar su nombre y su versión, y el editor comprueba que el YAML se pueda leer.
En ese código no se ha encontrado una validación contra los campos del modelo.
Qué comprueba el editor en la versión desplegada de la plataforma \citep{sphere} no se ha verificado.

Por eso la conformidad se establece por uso: una herramienta de este trabajo ingiere cada Datasheet en PRIME4API y la somete a diecisiete comprobaciones fijas y a una más por cada plan (Sección~\ref{sec:propuesta:prime4api}).
Cada propuesta declara así dos cosas separadas: en cuántas APIs del catálogo aparece la carencia y si hay traza de ejecución que muestre su efecto sobre el motor.
La primera es un argumento sobre el dominio, y la segunda, sobre el esquema.

Los elementos no soportados se reparten en dos clases, y cada una se prueba de forma distinta.
Una «carencia de expresividad» aparece cuando el proveedor publica algo que el esquema no tiene dónde poner.
Postmark vende seis volúmenes contratados para su cuota mensual, cada uno con su precio y su sobrecoste.
Una cuota de v0.3 solo tiene un \texttt{value} escalar y un \texttt{overage\_cost} único, así que no hay campo donde escribir la lista.
Esta carencia se detecta al modelar, porque el fichero no admite el dato.

Un «defecto de contrato» aparece cuando el esquema sí ofrece una forma, la Datasheet la usa bien y PRIME4API falla al leerla: se rompe el contrato implícito por el que el motor lee todo lo que la guía admite.
La guía admite dejar vacío un valor cuando no hay cifra que escribir.
En Postmark y en Mailchimp Transactional queda vacío el tope de destinatarios, que la fuente no pone, y en Resend el \textit{rate} del plan Enterprise, que existe y no se publica.
El motor convierte ese valor a número sin comprobar si está vacío, y la Datasheet, aun siendo conforme, deja sin responder diez comprobaciones en Postmark y en Mailchimp Transactional (Tabla~\ref{tab:verificacion}).
Este defecto solo se detecta al ejecutar el motor sobre el fichero, porque al modelar la Datasheet se escribe sin resistencia.

La primera consecuencia de la distinción está en la vía de promoción.
La recurrencia muestra que la carencia pertenece al dominio.
La traza muestra que el motor da un resultado erróneo o deja sin responder la operación, y esa prueba vale en una sola API para cualquiera de las dos clases.
Un defecto de contrato se descubre ejecutando el motor, así que llega ya con esa ejecución.
Una carencia de expresividad sin traza, como los volúmenes de Postmark, necesita aparecer en una segunda API.

La clase de carencia decide también dónde se arregla y en qué repositorio se comunica.
Los volúmenes de Postmark solo caben si cambia el esquema y la cuota admite una lista de volúmenes.
El campo vacío admite el arreglo en el esquema, en el motor o en los dos, y cada propuesta dice en cuál.
El cambio de la cuota se comunica en el repositorio de la guía, y la comprobación del campo vacío, en el del motor.
La separación se ha obtenido sobre un solo esquema, ACM2YAML v0.3, y un solo motor, PRIME4API, por lo que no se afirma que valga fuera de ese caso.
